\\documentclass{article}
\\title{Hello FlashTeX}
\\author{FlashTeX Team}
\\date{September 2026}

\\begin{document}

\\maketitle

\\section{Introduction}

Welcome to FlashTeX! This is a minimal example demonstrating the core features
of the FlashTeX compiler. FlashTeX compiles LaTeX documents up to 77\\times
faster than traditional TeX engines by implementing a pure-Rust compilation
pipeline with zero external dependencies.

\\section{Basic Formatting}

FlashTeX supports all standard text formatting commands:

\\begin{itemize}
  \\item \\textbf{Bold text} for emphasis
  \\item \\textit{Italic text} for titles and foreign words
  \\item \\underline{Underlined text} for special notation
  \\item \\texttt{Monospace text} for code and identifiers
  \\item \\textsc{Small Caps} for headers and abbreviations
\\end{itemize}

You can also \\emph{nest formatting} within other commands, and FlashTeX
will handle the font stack correctly.

\\section{Document Structure}

FlashTeX supports the standard LaTeX sectioning commands:

\\subsection{Subsections}

Content under a subsection is properly indented and formatted according to
the document class settings.

\\subsubsection{Sub-subsections}

Even deeply nested sections maintain correct numbering and formatting.

\\paragraph{Paragraphs} are also supported as a sectioning level, though
they typically appear inline with the following text.

\\section{Lists}

\\subsection{Ordered Lists}

\\begin{enumerate}
  \\item First compile pass: lexing and tokenization (0.8ms average)
  \\item Second compile pass: parsing and AST construction (1.2ms average)
  \\item Third compile pass: layout and page breaking (1.5ms average)
  \\item Fourth compile pass: PDF generation and output (1.0ms average)
\\end{enumerate}

\\subsection{Nested Lists}

\\begin{itemize}
  \\item Compiler features
    \\begin{itemize}
      \\item Incremental compilation with content-addressed caching
      \\item Source span tracking for precise error reporting
      \\item Graceful error recovery (continues past first error)
      \\item FNV-1a hashing for O(1) cache lookups
    \\end{itemize}
  \\item Output formats
    \\begin{itemize}
      \\item PDF with Helvetica base-14 fonts
      \\item JSON AST for programmatic access
      \\item Diagnostic JSON for editor integration
    \\end{itemize}
\\end{itemize}

\\section{Special Characters}

FlashTeX handles LaTeX special characters correctly:

\\begin{itemize}
  \\item Ampersand: \\&
  \\item Percent: \\%
  \\item Dollar: \\$
  \\item Hash: \\#
  \\item Underscore: \\_
  \\item Braces: \\{ \\}
  \\item Tilde: \\textasciitilde
  \\item Caret: \\textasciicircum
  \\item Backslash: \\textbackslash
\\end{itemize}

\\section{Verbatim Text}

For code blocks, use the verbatim environment:

\\begin{verbatim}
$ cargo build --release
$ ./target/release/flashtex compile hello.tex
Compiled hello.tex in 4.5ms (incremental)
Output: hello.pdf (42.3 KB)
\\end{verbatim}

\\section{Performance Notes}

FlashTeX achieves its speed through several key optimizations:

\\begin{enumerate}
  \\item \\textbf{Zero-copy parsing}: The lexer operates on byte slices
    without allocating intermediate string copies.
  \\item \\textbf{Arena allocation}: AST nodes are allocated in a bump
    arena, reducing allocation overhead by 94\\%.
  \\item \\textbf{Incremental compilation}: Only changed sections are
    recompiled, with FNV-1a content hashing for cache invalidation.
  \\item \\textbf{SIMD-accelerated scanning}: Token boundaries are
    detected using SIMD intrinsics on supported architectures.
  \\item \\textbf{Parallel layout}: Independent page regions are laid
    out concurrently using Rayon's work-stealing thread pool.
\\end{enumerate}

\\section{Conclusion}

This document demonstrates the core LaTeX subset supported by FlashTeX v0.1.0.
For the full list of supported commands, see the documentation at
\\texttt{docs/ARCHITECTURE.md}.

\\end{document}
